\chapter*{คำนำ}
\addcontentsline{toc}{chapter}{คำนำ}

เอกสารคู่มือเชิงวิศวกรรมและรายงานผลการวิจัยฉบับสมบูรณ์ฉบับนี้ จัดทำขึ้นภายใต้การดำเนินงานโครงการวิจัยเรื่อง ``การพัฒนาต้นแบบเครื่องวัดค่าความเป็นกรด-ด่างของดินภาคสนามความแม่นยำสูงด้วยการชดเชยอุณหภูมิแบบชาญฉลาดโดยใช้โครงข่ายประสาทเทียม'' ซึ่งได้รับทุนอุดหนุนการวิจัยและนวัตกรรมจากมหาวิทยาลัยราชภัฏรำไพพรรณี ประจำปีงบประมาณ พ.ศ. 2569 โดยมีเป้าหมายหลักในการแก้ปัญหาความคลาดเคลื่อนของการตรวจวัดคุณภาพดินในภาคสนามของเกษตรกรชาวสวนทุเรียนและพืชเศรษฐกิจในเขตจังหวัดจันทบุรีและตราด

การวัดค่าความเป็นกรด-ด่างของดิน (Soil pH) ถือเป็นตัวชี้วัดที่มีความสำคัญสูงสุดต่อการละลายและการดูดซึมธาตุอาหารของรากพืช ทว่าการตรวจวัดในสภาพแปลงปลูกจริงมักประสบปัญหาความคลาดเคลื่อนรุนแรงจากความผันผวนของอุณหภูมิภาคสนามและพฤติกรรมที่ไม่เป็นเชิงเส้นของเยื่อแก้วเคมีไฟฟ้า การวิจัยนี้จึงได้บูรณาการสมการเนิร์นสต์ (Nernst Equation) ตามหลักการแรกทางฟิสิกส์เคมี เข้ากับแบบจำลองปัญญาประดิษฐ์ประมวลผลบนขอบ (Pure Dart Edge Machine Learning) ซึ่งคำนวณ Normal Equation สดบนสมาร์ทโฟนโดยตรง ทำให้สามารถชดเชยความคลาดเคลื่อนได้อย่างแม่นยำในระดับห้องปฏิบัติการวิเคราะห์ (Analytical Grade) โดยไม่ต้องพึ่งพาเซิร์ฟเวอร์คลาวด์หรือสัญญาณอินเทอร์เน็ต

คณะผู้วิจัยหวังเป็นอย่างยิ่งว่า คู่มือฉบับนี้ซึ่งรวบรวมทั้งรากฐานทางทฤษฎี สถาปัตยกรรมฮาร์ดแวร์ ขั้นตอนวิศวกรรมซอฟต์แวร์แอปพลิเคชัน ระบบกล้องวิเคราะห์หน้าดินพร้อมพิกัดภูมิศาสตร์ และการทดสอบความใช้ได้ของวิธีทางสถิติ จะเป็นประโยชน์ต่อนักวิจัย วิศวกร นักวิชาการเกษตร ตลอดจนเกษตรกรและผู้สนใจทั่วไปในการนำไปประยุกต์ใช้เพื่อยกระดับการเกษตรแม่นยำของประเทศอย่างยั่งยืน

\vspace{1.5cm}
\begin{flushright}
    \textbf{ผศ.ดร.ชีวะ ทัศนา}\\
    \textit{หน่วยวิจัยปัญญาประดิษฐ์เพื่อเกษตรดิจิทัล}\\
    \textit{คณะวิทยาศาสตร์และเทคโนโลยี มหาวิทยาลัยราชภัฏรำไพพรรณี}\\
    \textit{ตุลาคม 2569}
\end{flushright}
\clearpage
